% !PRINTED 494
\weberpage{494}

%<<CONT
imprimitive numbers of the field. A symmetric function of the conjugate
numbers is a rational number. Among these symmetric functions two are of
especial importance, the sum and the product, of which the first is called
the \gesperrt{trace}, the second the \gesperrt{norm} of $\omega$. These two
numbers are denoted by $S(\omega)$ and $N(\omega)$.

Here, if among the conjugate numbers the same numbers occur several times,
these equal numbers are taken into the sum or the product as many times as
the degree of their frequency indicates.

Since in every such number field the four fundamental operations of
arithmetic can be carried out just as in the field of rational numbers, one
may also raise the question how far the fundamental arithmetical laws known
from the theory of rational numbers hold good in an arbitrary algebraic
number field. The first concern here is the decomposition of the integers
into their \gesperrt{prime factors}.

Since this decomposition does in general \gesperrt{not} succeed with the
numbers of the algebraic field themselves, an extension of the material of
computation is necessary in order to restore the simple laws, and such an
extension is possible in different senses. But these different extensions
must admit of being reduced to one another, or, more precisely stated, of
being placed in a unique relation to one another.

\gesperrt{Kummer} first solved the great problem, for the algebraic numbers
formed from roots of unity (the cyclotomic numbers), by the creation of the
\gesperrt{ideal numbers}\textsuperscript{1)}. Another entirely general
solution, subject to no exception, the problem has found through
\gesperrt{Dedekind}, who regarded the \gesperrt{ideals}\textsuperscript{2)} so named by him

\vspace{1em}\hrule\vspace{0.5em}

\noindent\textsuperscript{1)} \gesperrt{Kummer}, Theorie der idealen
Primfactoren der complexen Zahlen etc. \gesperrt{Crelle}'s Journal, Bd.\,35, 1846,
Bd.\,40, 1850. Abhandlungen der Berliner Akademie 1856. Sur la th\'eorie des
nombres complexes compos\'es de racines de l'unit\'e et de nombres entiers.
\gesperrt{Liouville}'s Journal, Bd.\,16, 1851.

\noindent\textsuperscript{2)} \gesperrt{Dedekind}, in the last Supplement of
the 2nd, 3rd and 4th editions of \gesperrt{Dirichlet's} Vorlesungen \"uber
Zahlentheorie (Braunschweig 1871, 1879, 1894). To be compared also: Sur la
th\'eorie des nombres entiers alg\'ebriques in the Bulletin of
\gesperrt{Darboux} and \gesperrt{Ho\"uel} (1\textsuperscript{re} s\'er. XI,
1877). Ueber den Zusammenhang zwischen der Theorie der Ideale und der
h\"oheren Congruenzen (Abhandlungen der Ges. der Wissensch. in G\"ottingen,
Bd.\,23,
%CONT>>
